% Фрагмент листочка для подключения через % Фрагмент листочка для подключения через % Фрагмент листочка для подключения через % Фрагмент листочка для подключения через \input{05.tex}.
% Черновик: условия можно редактировать; ответы и запасные задачи — в конце исходника.

\begin{center}
    \large\textbf{5. Отрицания. Рассуждения от противного}
\end{center}

\begin{enumerate}
    \item \textbf{Постройте отрицания утверждений.}


    \begin{enumerate}[label={},leftmargin=1.7em,itemsep=0.7em]
        \item[а)] \textit{Все взрослые пошли в кафе.} Это неверно. Значит, 

        \vspace{1.5cm}
        \item[б)] \textit{Этот пирожок с малиной или с черникой.} Это неверно. Значит, 

        \vspace{1.5cm}
        \item[в)] \textit{Завтра не будет математики и будет география.} Это неверно. Значит, 

        \vspace{1.5cm}


        \item[г)] \textit{В классе меньше 10 мальчиков.} Это неверно. Значит, 

        \vspace{1.5cm}


        \item[д)] \textit{Некоторые фуфырки не любят арригаторов.} Это неверно. Значит, 

        \vspace{1.5cm}
    \end{enumerate}
\end{enumerate}

\newpage

\begin{center}
\textbf{Метод доказательства от противного}
\end{center}

\textit{Дано:} утверждение А.

\textit{Нужно доказать:} утверждение Б.

\textit{Схема доказательства:}

1)  Предположим, что Б неверно. Сформулируем отрицание Б.

2) Выясним, что следует из отрицания Б. Сделаем логические переходы.

3) Получим противоречие с условием А или с известным фактом.

4) Сделаем вывод: наше предположение неверно, поэтому Б истинно. Что и требовалось доказать.

\begin{enumerate}[start=2,itemsep=0.8em]
    \item Докажите, что нельзя разложить 65 конфет на 11 кучек так, чтобы в любых двух кучках было разное количество конфет.

    \item Десять школьников на олимпиаде решили в сумме 35 задач. Среди них есть школьник, решивший ровно одну задачу, школьник, решивший ровно две задачи, и школьник, решивший ровно три задачи. Докажите, что хотя бы один школьник решил не менее пяти задач.

    \item На русско-французской встрече были только россияне и французы. Суммарное количество денег у французов оказалось больше, чем у россиян, а у женщин --- больше, чем у мужчин. Обязательно ли на встрече была француженка? Обоснуйте ответ.

    \item В вершинах куба расставлены числа 1, 2, 3, 4, 5, 6, 7, 8, каждое по одному разу. Докажите, что найдётся ребро, числа на концах которого отличаются не менее чем на 3.
\end{enumerate}

\newpage

\begin{center}
    \large\textbf{Самостоятельная работа. Вариант 1}
\end{center}

\textit{Автор решений:} \rule{6cm}{0.4pt}

\begin{enumerate}[itemsep=0.8em]
    \item Представьте дробь $\frac{7}{13}$ в виде десятичной дроби с периодом.

    \item Запишите отрицание утверждения: <<Каждый ученик решил хотя бы две задачи>>.

    \vspace{1.5cm}

    \item Найдите сотую цифру после запятой в десятичной записи числа $\dfrac{1}{7}$.

    \item В коробке лежит 10 карандашей. Cреди любых четырёх карандашей найдётся хотя бы один красный.

    Какие утверждения обязательно следуют из этого? (Их может быть несколько, одно или ни одного.)

    \begin{enumerate}[label={},leftmargin=1.7em,itemsep=0.4em]
        \item[а)] Карандашей не красного цвета не больше трёх.
        \item[б)] Красных карандашей не меньше семи.
        \item[в)] В коробке ровно семь красных карандашей.
        \item[г)] Среди любых пяти карандашей найдётся хотя бы два красных.
        \item[д)] Среди любых трёх карандашей найдётся хотя бы один красный.
    \end{enumerate}

    \newpage

    \item Не все ребята из шахматного кружка умеют плавать. Все участники соревнований по плаванию умеют плавать. Верно ли, что не все ребята из шахматного кружка участвуют в соревнованиях по плаванию? Объясните ответ.

    \vspace{5cm}

    \item  Числитель и знаменатель дроби — положительные числа. Числитель увеличили на 1, а знаменатель — на 100, при этом дробь увеличилась. Приведите пример такой дроби.
\end{enumerate}

\newpage

\begin{center}
    \large\textbf{Самостоятельная работа. Вариант 2}
\end{center}

\textit{Автор решений:} \rule{6cm}{0.4pt}

\begin{enumerate}[itemsep=0.8em]
    \item Представьте дробь $\frac{6}{13}$ в виде десятичной дроби с периодом.

    \item Запишите отрицание утверждения: <<Каждый школьник прочитал меньше двух книг>>.

    \vspace{1.5cm}

    \item Найдите сотую цифру после запятой в десятичной записи числа $\dfrac{2}{7}$.

    \item В коробке лежит 12 карандашей. Среди любых пяти карандашей найдётся хотя бы один синий.

    Какие утверждения обязательно следуют из этого? (Их может быть несколько, одно или ни одного.)

    \begin{enumerate}[label={},leftmargin=1.7em,itemsep=0.4em]
        \item[а)] Карандашей не синего цвета не больше четырёх.
        \item[б)] Синих карандашей не меньше восьми.
        \item[в)] В коробке ровно восемь синих карандашей.
        \item[г)] Среди любых шести карандашей найдётся хотя бы два синих.
        \item[д)] Среди любых четырёх карандашей найдётся хотя бы один синий.
    \end{enumerate}

    \newpage

    \item Не все ребята из математического кружка умеют кататься на коньках. Все участники соревнований по фигурному катанию умеют кататься на коньках. Верно ли, что не все ребята из математического кружка участвуют в соревнованиях по фигурному катанию? Объясните ответ.

    \vspace{5cm}

    \item Приведите пример двух обыкновенных дробей, разность которых в три раза больше их произведения. Приведите вычисления, обосновывающие это свойство.
\end{enumerate}
.
% Черновик: условия можно редактировать; ответы и запасные задачи — в конце исходника.

\begin{center}
    \large\textbf{5. Отрицания. Рассуждения от противного}
\end{center}

\begin{enumerate}
    \item \textbf{Постройте отрицания утверждений.}


    \begin{enumerate}[label={},leftmargin=1.7em,itemsep=0.7em]
        \item[а)] \textit{Все взрослые пошли в кафе.} Это неверно. Значит, 

        \vspace{1.5cm}
        \item[б)] \textit{Этот пирожок с малиной или с черникой.} Это неверно. Значит, 

        \vspace{1.5cm}
        \item[в)] \textit{Завтра не будет математики и будет география.} Это неверно. Значит, 

        \vspace{1.5cm}


        \item[г)] \textit{В классе меньше 10 мальчиков.} Это неверно. Значит, 

        \vspace{1.5cm}


        \item[д)] \textit{Некоторые фуфырки не любят арригаторов.} Это неверно. Значит, 

        \vspace{1.5cm}
    \end{enumerate}
\end{enumerate}

\newpage

\begin{center}
\textbf{Метод доказательства от противного}
\end{center}

\textit{Дано:} утверждение А.

\textit{Нужно доказать:} утверждение Б.

\textit{Схема доказательства:}

1)  Предположим, что Б неверно. Сформулируем отрицание Б.

2) Выясним, что следует из отрицания Б. Сделаем логические переходы.

3) Получим противоречие с условием А или с известным фактом.

4) Сделаем вывод: наше предположение неверно, поэтому Б истинно. Что и требовалось доказать.

\begin{enumerate}[start=2,itemsep=0.8em]
    \item Докажите, что нельзя разложить 65 конфет на 11 кучек так, чтобы в любых двух кучках было разное количество конфет.

    \item Десять школьников на олимпиаде решили в сумме 35 задач. Среди них есть школьник, решивший ровно одну задачу, школьник, решивший ровно две задачи, и школьник, решивший ровно три задачи. Докажите, что хотя бы один школьник решил не менее пяти задач.

    \item На русско-французской встрече были только россияне и французы. Суммарное количество денег у французов оказалось больше, чем у россиян, а у женщин --- больше, чем у мужчин. Обязательно ли на встрече была француженка? Обоснуйте ответ.

    \item В вершинах куба расставлены числа 1, 2, 3, 4, 5, 6, 7, 8, каждое по одному разу. Докажите, что найдётся ребро, числа на концах которого отличаются не менее чем на 3.
\end{enumerate}

\newpage

\begin{center}
    \large\textbf{Самостоятельная работа. Вариант 1}
\end{center}

\textit{Автор решений:} \rule{6cm}{0.4pt}

\begin{enumerate}[itemsep=0.8em]
    \item Представьте дробь $\frac{7}{13}$ в виде десятичной дроби с периодом.

    \item Запишите отрицание утверждения: <<Каждый ученик решил хотя бы две задачи>>.

    \vspace{1.5cm}

    \item Найдите сотую цифру после запятой в десятичной записи числа $\dfrac{1}{7}$.

    \item В коробке лежит 10 карандашей. Cреди любых четырёх карандашей найдётся хотя бы один красный.

    Какие утверждения обязательно следуют из этого? (Их может быть несколько, одно или ни одного.)

    \begin{enumerate}[label={},leftmargin=1.7em,itemsep=0.4em]
        \item[а)] Карандашей не красного цвета не больше трёх.
        \item[б)] Красных карандашей не меньше семи.
        \item[в)] В коробке ровно семь красных карандашей.
        \item[г)] Среди любых пяти карандашей найдётся хотя бы два красных.
        \item[д)] Среди любых трёх карандашей найдётся хотя бы один красный.
    \end{enumerate}

    \newpage

    \item Не все ребята из шахматного кружка умеют плавать. Все участники соревнований по плаванию умеют плавать. Верно ли, что не все ребята из шахматного кружка участвуют в соревнованиях по плаванию? Объясните ответ.

    \vspace{5cm}

    \item  Числитель и знаменатель дроби — положительные числа. Числитель увеличили на 1, а знаменатель — на 100, при этом дробь увеличилась. Приведите пример такой дроби.
\end{enumerate}

\newpage

\begin{center}
    \large\textbf{Самостоятельная работа. Вариант 2}
\end{center}

\textit{Автор решений:} \rule{6cm}{0.4pt}

\begin{enumerate}[itemsep=0.8em]
    \item Представьте дробь $\frac{6}{13}$ в виде десятичной дроби с периодом.

    \item Запишите отрицание утверждения: <<Каждый школьник прочитал меньше двух книг>>.

    \vspace{1.5cm}

    \item Найдите сотую цифру после запятой в десятичной записи числа $\dfrac{2}{7}$.

    \item В коробке лежит 12 карандашей. Среди любых пяти карандашей найдётся хотя бы один синий.

    Какие утверждения обязательно следуют из этого? (Их может быть несколько, одно или ни одного.)

    \begin{enumerate}[label={},leftmargin=1.7em,itemsep=0.4em]
        \item[а)] Карандашей не синего цвета не больше четырёх.
        \item[б)] Синих карандашей не меньше восьми.
        \item[в)] В коробке ровно восемь синих карандашей.
        \item[г)] Среди любых шести карандашей найдётся хотя бы два синих.
        \item[д)] Среди любых четырёх карандашей найдётся хотя бы один синий.
    \end{enumerate}

    \newpage

    \item Не все ребята из математического кружка умеют кататься на коньках. Все участники соревнований по фигурному катанию умеют кататься на коньках. Верно ли, что не все ребята из математического кружка участвуют в соревнованиях по фигурному катанию? Объясните ответ.

    \vspace{5cm}

    \item Приведите пример двух обыкновенных дробей, разность которых в три раза больше их произведения. Приведите вычисления, обосновывающие это свойство.
\end{enumerate}
.
% Черновик: условия можно редактировать; ответы и запасные задачи — в конце исходника.

\begin{center}
    \large\textbf{5. Отрицания. Рассуждения от противного}
\end{center}

\begin{enumerate}
    \item \textbf{Постройте отрицания утверждений.}


    \begin{enumerate}[label={},leftmargin=1.7em,itemsep=0.7em]
        \item[а)] \textit{Все взрослые пошли в кафе.} Это неверно. Значит, 

        \vspace{1.5cm}
        \item[б)] \textit{Этот пирожок с малиной или с черникой.} Это неверно. Значит, 

        \vspace{1.5cm}
        \item[в)] \textit{Завтра не будет математики и будет география.} Это неверно. Значит, 

        \vspace{1.5cm}


        \item[г)] \textit{В классе меньше 10 мальчиков.} Это неверно. Значит, 

        \vspace{1.5cm}


        \item[д)] \textit{Некоторые фуфырки не любят арригаторов.} Это неверно. Значит, 

        \vspace{1.5cm}
    \end{enumerate}
\end{enumerate}

\newpage

\begin{center}
\textbf{Метод доказательства от противного}
\end{center}

\textit{Дано:} утверждение А.

\textit{Нужно доказать:} утверждение Б.

\textit{Схема доказательства:}

1)  Предположим, что Б неверно. Сформулируем отрицание Б.

2) Выясним, что следует из отрицания Б. Сделаем логические переходы.

3) Получим противоречие с условием А или с известным фактом.

4) Сделаем вывод: наше предположение неверно, поэтому Б истинно. Что и требовалось доказать.

\begin{enumerate}[start=2,itemsep=0.8em]
    \item Докажите, что нельзя разложить 65 конфет на 11 кучек так, чтобы в любых двух кучках было разное количество конфет.

    \item Десять школьников на олимпиаде решили в сумме 35 задач. Среди них есть школьник, решивший ровно одну задачу, школьник, решивший ровно две задачи, и школьник, решивший ровно три задачи. Докажите, что хотя бы один школьник решил не менее пяти задач.

    \item На русско-французской встрече были только россияне и французы. Суммарное количество денег у французов оказалось больше, чем у россиян, а у женщин --- больше, чем у мужчин. Обязательно ли на встрече была француженка? Обоснуйте ответ.

    \item В вершинах куба расставлены числа 1, 2, 3, 4, 5, 6, 7, 8, каждое по одному разу. Докажите, что найдётся ребро, числа на концах которого отличаются не менее чем на 3.
\end{enumerate}

\newpage

\begin{center}
    \large\textbf{Самостоятельная работа. Вариант 1}
\end{center}

\textit{Автор решений:} \rule{6cm}{0.4pt}

\begin{enumerate}[itemsep=0.8em]
    \item Представьте дробь $\frac{7}{13}$ в виде десятичной дроби с периодом.

    \item Запишите отрицание утверждения: <<Каждый ученик решил хотя бы две задачи>>.

    \vspace{1.5cm}

    \item Найдите сотую цифру после запятой в десятичной записи числа $\dfrac{1}{7}$.

    \item В коробке лежит 10 карандашей. Cреди любых четырёх карандашей найдётся хотя бы один красный.

    Какие утверждения обязательно следуют из этого? (Их может быть несколько, одно или ни одного.)

    \begin{enumerate}[label={},leftmargin=1.7em,itemsep=0.4em]
        \item[а)] Карандашей не красного цвета не больше трёх.
        \item[б)] Красных карандашей не меньше семи.
        \item[в)] В коробке ровно семь красных карандашей.
        \item[г)] Среди любых пяти карандашей найдётся хотя бы два красных.
        \item[д)] Среди любых трёх карандашей найдётся хотя бы один красный.
    \end{enumerate}

    \newpage

    \item Не все ребята из шахматного кружка умеют плавать. Все участники соревнований по плаванию умеют плавать. Верно ли, что не все ребята из шахматного кружка участвуют в соревнованиях по плаванию? Объясните ответ.

    \vspace{5cm}

    \item  Числитель и знаменатель дроби — положительные числа. Числитель увеличили на 1, а знаменатель — на 100, при этом дробь увеличилась. Приведите пример такой дроби.
\end{enumerate}

\newpage

\begin{center}
    \large\textbf{Самостоятельная работа. Вариант 2}
\end{center}

\textit{Автор решений:} \rule{6cm}{0.4pt}

\begin{enumerate}[itemsep=0.8em]
    \item Представьте дробь $\frac{6}{13}$ в виде десятичной дроби с периодом.

    \item Запишите отрицание утверждения: <<Каждый школьник прочитал меньше двух книг>>.

    \vspace{1.5cm}

    \item Найдите сотую цифру после запятой в десятичной записи числа $\dfrac{2}{7}$.

    \item В коробке лежит 12 карандашей. Среди любых пяти карандашей найдётся хотя бы один синий.

    Какие утверждения обязательно следуют из этого? (Их может быть несколько, одно или ни одного.)

    \begin{enumerate}[label={},leftmargin=1.7em,itemsep=0.4em]
        \item[а)] Карандашей не синего цвета не больше четырёх.
        \item[б)] Синих карандашей не меньше восьми.
        \item[в)] В коробке ровно восемь синих карандашей.
        \item[г)] Среди любых шести карандашей найдётся хотя бы два синих.
        \item[д)] Среди любых четырёх карандашей найдётся хотя бы один синий.
    \end{enumerate}

    \newpage

    \item Не все ребята из математического кружка умеют кататься на коньках. Все участники соревнований по фигурному катанию умеют кататься на коньках. Верно ли, что не все ребята из математического кружка участвуют в соревнованиях по фигурному катанию? Объясните ответ.

    \vspace{5cm}

    \item Приведите пример двух обыкновенных дробей, разность которых в три раза больше их произведения. Приведите вычисления, обосновывающие это свойство.
\end{enumerate}
.
% Черновик: условия можно редактировать; ответы и запасные задачи — в конце исходника.

\begin{center}
    \large\textbf{5. Отрицания. Рассуждения от противного}
\end{center}

\begin{enumerate}
    \item \textbf{Постройте отрицания утверждений.}


    \begin{enumerate}[label={},leftmargin=1.7em,itemsep=0.7em]
        \item[а)] \textit{Все взрослые пошли в кафе.} Это неверно. Значит, 

        \vspace{1.5cm}
        \item[б)] \textit{Этот пирожок с малиной или с черникой.} Это неверно. Значит, 

        \vspace{1.5cm}
        \item[в)] \textit{Завтра не будет математики и будет география.} Это неверно. Значит, 

        \vspace{1.5cm}


        \item[г)] \textit{В классе меньше 10 мальчиков.} Это неверно. Значит, 

        \vspace{1.5cm}


        \item[д)] \textit{Некоторые фуфырки не любят арригаторов.} Это неверно. Значит, 

        \vspace{1.5cm}
    \end{enumerate}
\end{enumerate}

\newpage

\begin{center}
\textbf{Метод доказательства от противного}
\end{center}

\textit{Дано:} утверждение А.

\textit{Нужно доказать:} утверждение Б.

\textit{Схема доказательства:}

1)  Предположим, что Б неверно. Сформулируем отрицание Б.

2) Выясним, что следует из отрицания Б. Сделаем логические переходы.

3) Получим противоречие с условием А или с известным фактом.

4) Сделаем вывод: наше предположение неверно, поэтому Б истинно. Что и требовалось доказать.

\begin{enumerate}[start=2,itemsep=0.8em]
    \item Докажите, что нельзя разложить 65 конфет на 11 кучек так, чтобы в любых двух кучках было разное количество конфет.

    \item Десять школьников на олимпиаде решили в сумме 35 задач. Среди них есть школьник, решивший ровно одну задачу, школьник, решивший ровно две задачи, и школьник, решивший ровно три задачи. Докажите, что хотя бы один школьник решил не менее пяти задач.

    \item На русско-французской встрече были только россияне и французы. Суммарное количество денег у французов оказалось больше, чем у россиян, а у женщин --- больше, чем у мужчин. Обязательно ли на встрече была француженка? Обоснуйте ответ.

    \item В вершинах куба расставлены числа 1, 2, 3, 4, 5, 6, 7, 8, каждое по одному разу. Докажите, что найдётся ребро, числа на концах которого отличаются не менее чем на 3.
\end{enumerate}

\newpage

\begin{center}
    \large\textbf{Самостоятельная работа. Вариант 1}
\end{center}

\textit{Автор решений:} \rule{6cm}{0.4pt}

\begin{enumerate}[itemsep=0.8em]
    \item Представьте дробь $\frac{7}{13}$ в виде десятичной дроби с периодом.

    \item Запишите отрицание утверждения: <<Каждый ученик решил хотя бы две задачи>>.

    \vspace{1.5cm}

    \item Найдите сотую цифру после запятой в десятичной записи числа $\dfrac{1}{7}$.

    \item В коробке лежит 10 карандашей. Cреди любых четырёх карандашей найдётся хотя бы один красный.

    Какие утверждения обязательно следуют из этого? (Их может быть несколько, одно или ни одного.)

    \begin{enumerate}[label={},leftmargin=1.7em,itemsep=0.4em]
        \item[а)] Карандашей не красного цвета не больше трёх.
        \item[б)] Красных карандашей не меньше семи.
        \item[в)] В коробке ровно семь красных карандашей.
        \item[г)] Среди любых пяти карандашей найдётся хотя бы два красных.
        \item[д)] Среди любых трёх карандашей найдётся хотя бы один красный.
    \end{enumerate}

    \newpage

    \item Не все ребята из шахматного кружка умеют плавать. Все участники соревнований по плаванию умеют плавать. Верно ли, что не все ребята из шахматного кружка участвуют в соревнованиях по плаванию? Объясните ответ.

    \vspace{5cm}

    \item  Числитель и знаменатель дроби — положительные числа. Числитель увеличили на 1, а знаменатель — на 100, при этом дробь увеличилась. Приведите пример такой дроби.
\end{enumerate}

\newpage

\begin{center}
    \large\textbf{Самостоятельная работа. Вариант 2}
\end{center}

\textit{Автор решений:} \rule{6cm}{0.4pt}

\begin{enumerate}[itemsep=0.8em]
    \item Представьте дробь $\frac{6}{13}$ в виде десятичной дроби с периодом.

    \item Запишите отрицание утверждения: <<Каждый школьник прочитал меньше двух книг>>.

    \vspace{1.5cm}

    \item Найдите сотую цифру после запятой в десятичной записи числа $\dfrac{2}{7}$.

    \item В коробке лежит 12 карандашей. Среди любых пяти карандашей найдётся хотя бы один синий.

    Какие утверждения обязательно следуют из этого? (Их может быть несколько, одно или ни одного.)

    \begin{enumerate}[label={},leftmargin=1.7em,itemsep=0.4em]
        \item[а)] Карандашей не синего цвета не больше четырёх.
        \item[б)] Синих карандашей не меньше восьми.
        \item[в)] В коробке ровно восемь синих карандашей.
        \item[г)] Среди любых шести карандашей найдётся хотя бы два синих.
        \item[д)] Среди любых четырёх карандашей найдётся хотя бы один синий.
    \end{enumerate}

    \newpage

    \item Не все ребята из математического кружка умеют кататься на коньках. Все участники соревнований по фигурному катанию умеют кататься на коньках. Верно ли, что не все ребята из математического кружка участвуют в соревнованиях по фигурному катанию? Объясните ответ.

    \vspace{5cm}

    \item Приведите пример двух обыкновенных дробей, разность которых в три раза больше их произведения. Приведите вычисления, обосновывающие это свойство.
\end{enumerate}
